% 原书 PDF 274，印刷页 262：从9.4标题开始；此前习题9.3.14由9.3节文件负责。
\section{Picard 概形}
\label{ch09s04:section}
\begingroup
\def\PS{\mathbf{Pic}_{X/S}}
\def\PE{\Pic_{(X/S)(\text{\'et})}}
\def\PF{\Pic_{(X/S)(\mathrm{fppf})}}
\def\L{\mathcal L}\def\N{\mathcal N}\def\Q{\mathcal Q}\def\PP{\mathcal P}
本节证明 Grothendieck 关于 Picard 概形的主定理：若 $X/S$ 射影且平坦，几何纤维为整概形，则 Picard 概形存在；事实上，函子 $\PE$ 可表。证明采用 Grothendieck 利用函子规定拼接的方法。基本理论见 [EGAG，第 0 章，第 4.5 节，第 102--107 页]，并在 [EGAG，第 1 章，第 9 节，第 354--401 页] 中用于构造 Grassmann 簇及相关的参数概形。这里的构造也运用基本理论，但更精细、更复杂，因为它不只使用 Zariski 拓扑，还使用平展拓扑。

\begin{origstatement}{定义}{9.4.1}\label{ch09s04:9.4.1}
若定义 9.2.2 中四个相对 Picard 函子中的任一个可表，则称表示它的概形为 \emph{Picard 概形}，记为 $\PS$。此外，简单地说，Picard 概形 $\PS$ \emph{存在}。
\end{origstatement}
注意，尽管有四个相对 Picard 函子，Picard 概形至多只有一个。当然，若某个函子可表，则表示概形唯一确定到唯一同构，该同构保持给定函子与表示概形的点函子之间的等同。但是，这里还有更多内容。

事实上，例如假定 $\PE$ 由 $\PS$ 表示。则由下降理论，$\PE$ 已经是 fppf 拓扑中的层；所以它等于其相关联层 $\PF$。因此，$\PF$ 也由 $\PS$ 表示。另一方面，$\Pic_{X/S}$ 可能可表，也可能不可表；但若可表，它也必由 $\PS$ 表示。

\begin{origstatement}{习题}{9.4.2}\label{ch09s04:9.4.2}
假定 $\PS$ 存在，且 $\mathcal O_S\xrightarrow{\sim}f_*\mathcal O_X$ 普遍成立。设 $T$ 为 $S$ 概形，$\L$ 为 $X_T$ 上可逆层。证明，存在子概形 $N\subset T$ 及 $N$ 上的可逆层 $\N$，具有以下三个性质：第一，$\L_N\simeq f_N^*\N$；第二，给定任意 $S$ 映射 $t:T'\to T$，若对 $T'$ 上某个可逆层 $\N'$ 有 $\L_{T'}\simeq f_{T'}^*\N'$，则 $t$ 必通过 $N$ 分解，且 $\N'\simeq t^*\N$；第三，若 $\PS$ 分离，则 $N$ 为闭子概形。还要证明，前两个性质唯一确定 $N$，并在同构意义下确定 $\N$。
\end{origstatement}
\begin{origstatement}{习题}{9.4.3}\label{ch09s04:9.4.3}
假定 $\PS$ 存在。$X\times\PS$ 上的可逆层 $\PP$ 称为\emph{万有层}（universal sheaf）或 \emph{Poincaré 层}，若 $\PP$ 具有以下性质：给定任意 $S$ 概形 $T$ 以及 $X_T$ 上任意可逆层 $\L$，存在唯一 $S$ 映射 $h:T\to\PS$，使得对 $T$ 上某个可逆层 $\N$，有
\[
\L\simeq(1\times h)^*\PP\otimes f_T^*\N.
\]
证明，万有层 $\PP$ 存在，当且仅当 $\PS$ 表示 $\Pic_{X/S}$。
% 原书 PDF 275，印刷页 263；习题9.4.3续。

假定 $\mathcal O_S\xrightarrow{\sim}f_*\mathcal O_X$ 普遍成立。证明，若 $\PP$ 存在，则它在与 $\PS$ 上唯一一个可逆层的拉回作张量积的意义下唯一。

证明，若 $f$ 还有一个截面，则万有层 $\PP$ 存在。

找出一个不存在万有层 $\PP$ 的例子。
\end{origstatement}
\begin{origstatement}{习题}{9.4.4}\label{ch09s04:9.4.4}
假定 $\PS$ 存在，并设 $S'$ 为 $S$ 概形。证明，$\mathbf{Pic}_{X_{S'}/S'}$ 也存在，且事实上
\[
\mathbf{Pic}_{X_{S'}/S'}=\PS\times_S S'.
\]
因此，Picard 概形的形成与基变换交换。

找出一个例子，使 $\mathbf{Pic}_{X_{S'}/S'}$ 表示 $\Pic_{X_{S'}/S'}$，而 $\PS$ 不表示 $\Pic_{X/S}$。
\end{origstatement}
\begin{origstatement}{习题}{9.4.5}\label{ch09s04:9.4.5}
假定 $\PS$ 存在，且它表示 $\PF$，或 $\mathcal O_S\xrightarrow{\sim}f_*\mathcal O_X$ 普遍成立。证明，$\PS$ 的概形点与 $X/S$ 的纤维上的可逆层 $\L$ 的类自然一一对应。按定义，一个类由某个 $X_k$ 上的 $\L$ 表示，其中 $k$ 是包含某个（概形）点 $s\in S$ 的剩余域 $k_s$ 的域；$X_{k'}$ 上的 $\L'$ 表示同一个类，当且仅当存在包含前两个域的第三个域 $k''$，使 $\L|X_{k''}\simeq\L'|X_{k''}$。
\end{origstatement}
\begin{origstatement}{定义}{9.4.6}\label{ch09s04:9.4.6}
\emph{Abel 映射}是函子的自然映射
\[
A_{X/S}(T):\Div_{X/S}(T)\longrightarrow\Pic_{X/S}(T),
\]
定义为将 $X_T/T$ 上的相对有效除子 $D$ 映到层 $\mathcal O_{X_T}(D)$。目标 $\Pic_{X/S}$ 可以换成它的任一相关联层。若 $\PS$ 存在，则术语 Abel 映射也可指相应的概形映射
\[
\mathbf A_{X/S}:\mathbf{Div}_{X/S}\longrightarrow\PS.
\]
\end{origstatement}
\begin{origstatement}{习题}{9.4.7}\label{ch09s04:9.4.7}
假定 $X/S$ 固有且平坦，几何纤维为整概形。假定 $\PS$ 存在，以 $P$ 表示它。经 Abel 映射，将 $\mathbf{Div}_{X/S}$ 视为 $P$ 概形。假定存在万有层 $\PP$，并令 $\Q$ 为如第 9.3.10 小节中那样与 $\PP$ 配备的 $P$ 上的层。证明，作为 $P$ 概形，有 $\mathbb P(\Q)=\mathbf{Div}_{X/S}$。
\end{origstatement}
\begin{origstatement}{定理}{9.4.8}\label{ch09s04:9.4.8}
（主定理）假定 $f:X\to S$ 在 $S$ 上 Zariski 局部地为射影态射，并且平坦，几何纤维为整概形。

(1) 则 $\PS$ 存在，在 $S$ 上分离且局部有限型，并表示 $\PE$。

(2) 若 $S$ 还为 Noether 概形，且 $X/S$ 射影，则 $\PS$ 为开子概形的不交并，其中每个开子概形又是递增的拟射影 $S$ 开子概形的并。
\end{origstatement}
\begin{proof}
由 [EGAG，(0，4.5.5)，第 106 页]，表示 $S$ 概形范畴上的一个 Zariski 层，是 $S$ 上的局部问题。此外，证明 $S$ 概形分离且局部有限型，也是 $S$ 上的局部问题。因此，为证明 (1)，可以假定 $S$ 为 Noether 概形，且 $X/S$ 射影。

显然，若一个 $S$ 概形是分离开子概形的不交并，或是分离开子概形的递增并，则它分离。因此 (1) 由 (2) 得到。

为证明 (2)，由 Yoneda 引理，可以将概形 $T$ 与其点函子等同，从而将概形范畴看作函子范畴的满子范畴。为减轻记号，也以 $T$ 表示这一函子。
% 原书 PDF 276，印刷页 264。
并且，除了说函子可表之外，也说这个函子是概形。此外，置
\[
P:=\PE.
\]
注意 $P(T)=\Hom(T,P)$。

给定多项式 $\varphi\in\mathbb Q[n]$，令 $P^\varphi\subset P$ 为下列预层的相关联平展子层：该预层的 $T$ 点由 $X_T$ 上的可逆层 $\L$ 表示，且满足
\begin{equation}
\chi(X_t,\L_t^{-1}(n))=\varphi(n)\qquad\text{对所有 }t\in T.\tag{9.4.8.1}
\end{equation}
注意，这个预层定义良好，因为在任意基变换 $p:T'\to T$ 后，(9.4.8.1) 仍然成立；事实上，对任意 $t'\in T'$、任意 $i$ 和任意 $n$，有
\[
H^i(X_{t'},\L_{t'}^{-1}(n))
=H^i(X_{p(t')},\L_{p(t')}^{-1}(n))\otimes_{k_t}k_{t'},
\]
因为由 [Har77，命题 III，9.3，第 255 页]，上同调与平坦基变换交换。因此，$P^\varphi$ 也定义良好。

固定映射 $T\to P$，以平展覆盖 $p:T'\to T$ 及 $X_{T'}$ 上的可逆层 $\L'$ 表示它。考虑如下定义的子集 $T'^\varphi\subset T'$：
\[
T'^\varphi:=\{t'\in T'\mid\chi(X_{t'},\L_{t'}'{}^{-1}(n))=\varphi(n)\}.
\]
则由 [EGAIII2，7.9.11]，$T'^\varphi$ 为开集。

置 $T^\varphi:=p(T'^\varphi)$。则 $T^\varphi\subset T$ 为开集，因为 $T'^\varphi\subset T'$ 为开集且 $p$ 平展。此外，$T'^\varphi=p^{-1}(T^\varphi)$。事实上，设 $t'\in p^{-1}(T^\varphi)$，记 $p(t')=p(t'_1)$，其中 $t'_1\in T'^\varphi$。现在，存在平展覆盖 $T''\to T'\times_TT'$，使 $\L'$ 到 $X_{T''}$ 的两个拉回同构。取 $t''\in T''$，使其在第一个映射 $T''\to T'$ 下的像为 $t'\in T'$，在第二个映射下的像为 $t'_1\in T'$。则
\[
\chi(X_{t'},\L_{t'}'{}^{-1}(n))
=\chi(X_{t''},\L_{t''}'{}^{-1}(n))
=\chi(X_{t'_1},\L_{t'_1}'{}^{-1}(n))=\varphi(n).
\]
所以 $t'\in T'^\varphi$。因而 $T'^\varphi\supset p^{-1}(T^\varphi)$。故 $T'^\varphi=p^{-1}(T^\varphi)$。

此外，$T^\varphi$ 是（表示）函子的纤维积 $P^\varphi\times_P T$。事实上，为验证它们有相同的 $R$ 点，设 $r:R\to T$ 为映射；形成 $R':=R\times_TT'$ 和 $r':R'\to T'$。假定 $r$ 通过 $T^\varphi$ 分解。则 $r'$ 通过 $T'^\varphi$ 分解。所以，基本上由定义，$R'\to T'\to P$ 通过 $P^\varphi$ 分解。现在，$R'\to R$ 为平展覆盖。因此，由于 $P^\varphi$ 是平展层，$R\to T\to P$ 通过 $P^\varphi$ 分解。

反之，假定 $R\to T\to P$ 通过 $P^\varphi$ 分解。则 $R\to P$ 由平展覆盖 $R''\to R$ 及 $X_{R''}$ 上的可逆层 $\L''$ 定义，且对所有 $u\in R''$，有 $\chi(X_u,\L_u''{}^{-1}(n))=\varphi(n)$。由于 $\L''$ 与 $\L'$ 均定义 $R\to P$，存在平展覆盖 $R'''\to R''\times_RR'$，使 $\L''$ 和 $\L'$ 到 $X_{R'''}$ 的拉回同构。因此，$r':R'\to T'$ 的像位于 $T'^\varphi$ 中。但后者为开集，所以 $r'$ 通过它分解。因此，$r:R\to T$ 通过 $T^\varphi$ 分解。故 $T^\varphi$ 和 $P^\varphi\times_P T$ 有相同的 $R$ 点。

令 $\varphi$ 变动。显然，各 $T'^\varphi$ 互不相交，并覆盖 $T'$。所以各 $T^\varphi$ 互不相交，并覆盖 $T$。因此，由一般结果 [EGAG，(0，4.5.4)，第 103 页]，若各 $P^\varphi$ 是（可由）概形（表示的函子），则 $P$ 为它们的不交并。于是，只需将每个 $P^\varphi$ 表示成拟射影 $S$ 开子概形的递增并。

固定 $\varphi$。给定 $m\in\mathbb Z$，令 $P_m^\varphi\subset P^\varphi$ 为下列预层的相关联平展子层：其 $T$ 点由 $X_T$ 上的 $\L$ 表示，除 (9.4.8.1) 外，还要求
\begin{equation}
R^if_{T*}\L(n)=0\qquad\text{对所有 }i\geq1\text{ 及 }n\geq m.\tag{9.4.8.2}
\end{equation}
% 原书 PDF 277，印刷页 265。
注意，这个预层定义良好，因为在任意基变换 $p:T'\to T$ 后，(9.4.8.2) 仍然成立；下面说明这一点。

首先，说明 (9.4.8.2) 等价于下列条件：
\begin{equation}
H^i(X_t,\L_t(n))=0\quad\text{对所有 }i\geq1,\ n\geq m\text{ 及 }t\in T.\tag{9.4.8.3}
\end{equation}
事实上，(9.4.8.3) 蕴含 (9.4.8.2)；给定任意 $i,t,n$，若 $H^i(X_t,\L_t(n))=0$，则由 [EGAIII2，7.5.3] 或 [OB72，推论 2.1，第 68 页]，对 $T$ 上所有拟凝聚层 $\N$，有 $R^if_{T*}(\L(n)\otimes f_T^*\N)_t=0$。

反之，假定 (9.4.8.2)。固定 $t,n$。对 $i$ 作降归纳，证明 $H^i(X_t,\L_t(n))$ 消失。由 Serre 定理 [EGAIII1，2.2.2]，当 $i\gg1$ 时它消失。假定对某个 $i\geq2$ 它消失。如刚才所述，则 $R^if_{T*}(\L(n)\otimes f_T^*\N)_t$ 对 $T$ 上所有拟凝聚层 $\N$ 都消失。因此，由高阶正像的长正合列，$R^{i-1}f_{T*}(\L(n)\otimes f_T^*\N)_t$ 关于 $\N$ 右正合。所以，根据一般原理，有函子的自然同构
\[
R^{i-1}f_{T*}(\L(n))_t\otimes\N_t
\xrightarrow{\sim}R^{i-1}f_{T*}(\L(n)\otimes f_T^*\N)_t.
\]
由于 (9.4.8.2) 成立，定义域和目标均消失。取 $\N:=k_t$，便得 $H^{i-1}(X_t,\L_t(n))$ 的消失。因此，(9.4.8.2) 蕴含 (9.4.8.3)。

最后，对任意 $t'\in T'$、任意 $i$ 及任意 $n$，有
\[
H^i(X_{t'},\L_{t'}(n))
=H^i(X_{p(t')},\L_{p(t')}(n))\otimes_{k_t}k_{t'},
\]
因为上同调与平坦基变换交换。因此，基变换 $p:T'\to T$ 后，(9.4.8.3) 仍然成立；从而 (9.4.8.2) 也仍成立。所以预层定义良好，$P_m^\varphi$ 也定义良好。

与考察 $P^\varphi\times_P T$ 时基本相同的论证表明：给定映射 $T\to P^\varphi$，当 $m$ 变动时，积 $P_m^\varphi\times_{P^\varphi}T$ 构成 $T$ 的开子概形的嵌套序列，其并为 $T$。论证中的关键变化在于证明开性。此时不使用 [EGAIII2，7.9.11]，而使用 Serre 定理 [EGAIII1，2.2.2] 的下列部分：给定 Noether 环 $A$ 上的射影概形上的凝聚层 $\mathcal F$，仅有有限多个 $i\geq1$、$n\geq m$，使得 $H^i(\mathcal F(n))$ 非零，且这些非零 $A$ 模都有限生成。因此，若存在 $A$ 的素理想 $\mathfrak p$，使得对所有 $i\geq1$、$n\geq m$ 均有 $H^i(\mathcal F(n))_{\mathfrak p}=0$，则存在 $a\notin\mathfrak p$，使得对所有 $i\geq1$、$n\geq m$ 均有 $H^i(\mathcal F(n))_a=0$。

因此，再次由 [EGAG，(0，4.5.4)，第 103 页]，只需以拟射影 $S$ 概形表示每个 $P_m^\varphi$。

固定 $\varphi,m$。置 $\varphi_0(n):=\varphi(m+n)$。则有函子同构 $P_m^\varphi\xrightarrow{\sim}P_0^{\varphi_0}$，定义如下。首先，定义 $\Pic_{X/S}$ 的自态射 $\varepsilon$，将 $X_T$ 上的可逆层 $\L$ 映到 $\L(m)$。显然，$\varepsilon$ 是自同构。所以 $\varepsilon$ 诱导相关联层 $P$ 的自同构 $\varepsilon^+$。显然，$\varepsilon^+$ 将 $P_m^\varphi$ 映到 $P_0^{\varphi_0}$ 上。因此，只需以拟射影 $S$ 概形表示 $P_0^{\varphi_0}$。

由 [EGAIII2，7.9.11]，函数 $s\mapsto\chi(X_s,\mathcal O_{X_s}(n))$ 在 $S$ 上局部常值。因此，将 $S$ 换成一个既开又闭的子集，可以假定它为常值。置 $\psi(n):=\chi(X_s,\mathcal O_{X_s}(n))$。

考虑 Abel 映射 $A_{X/S}:\mathbf{Div}_{X/S}\to P$。注意，由定理 9.3.7，$\mathbf{Div}_{X/S}$ 是概形，事实上是 Hilbert 概形 $\mathbf{Hilb}_{X/S}$ 的开子概形。形成积 $P_0^{\varphi_0}\times_P\mathbf{Div}_{X/S}$。根据上面已证明的内容，它是概形 $Z$，事实上是 $\mathbf{Div}_{X/S}$ 的开子概形。置 $\theta(n):=\psi(n)-\varphi_0(n)$。显然，$Z$ 位于 $\mathbf{Hilb}_{X/S}^{\theta(n)}$ 中，后者在 $S$ 上射影。因此 $Z$ 拟射影。
% 原书 PDF 278，印刷页 266。

现在证明投影 $\alpha:Z\to P_0^{\varphi_0}$ 为平展层的满射。换言之，给定 $T$ 及 $\lambda\in P_0^{\varphi_0}(T)$，需要找出平展覆盖 $T_1\to T$ 以及 $\lambda_1\in Z(T_1)$，使 $\alpha(\lambda_1)\in P_0^{\varphi_0}(T_1)$ 等于 $\lambda$ 的像。

以平展覆盖 $p:T'\to T$ 及 $X_{T'}$ 上的可逆层 $\L'$ 表示 $\lambda$。几乎直接由定义，积 $T'\times_{P_0^{\varphi_0}}Z$ 等于 $\operatorname{LinSys}_{\L'/X_{T'}/T'}$。所以，由定理 9.3.13，这个积等于 $\mathbb P(\Q)$，其中 $\Q$ 是如第 9.3.10 小节中那样与 $\L'$ 配备的 $\mathcal O_{T'}$ 模。现在 $m=0$，因此由 (9.4.8.3)，对所有 $t\in T'$，有 $H^1(X_t,\L'_t)=0$。由于第 9.3.10 小节中的 (v) 蕴含 (i)，$\Q$ 局部自由。

因此，$\mathbb P(\Q)$ 在 $T'$ 上光滑。所以由 [EGAIV4，17.16.3(ii)]，存在平展覆盖 $T_1\to T'$ 及 $T'$ 映射 $T_1\to\mathbb P(\Q)$。于是，复合 $T_1\to\mathbb P(\Q)\to Z\to P_0^{\varphi_0}$ 等于复合 $T_1\to T'\to T\to P_0^{\varphi_0}$。换言之，映射 $T_1\to Z$ 是 $\lambda_1\in Z(T_1)$，使 $\alpha(\lambda_1)\in P_0^{\varphi_0}(T_1)$ 等于 $\lambda$ 的像。由于复合 $T_1\to T'\to T$ 为平展覆盖，$\alpha$ 因而为平展层的满射。

显然，映射 $\alpha:Z\to P_0^{\varphi_0}$ 由 $X_Z/Z$ 上万有相对有效除子所配备的可逆层定义。因此，在上面取 $T:=Z$、$T':=T$，可知积 $Z\times_{P_0^{\varphi_0}}Z$ 是概形，且第一投影光滑且射影。因此，定理由下列一般引理得到。
\end{proof}
\begin{origstatement}{引理}{9.4.9}\label{ch09s04:9.4.9}
设 $\alpha:Z\to P$ 为平展层的映射，并置 $R:=Z\times_P Z$。假定 $\alpha$ 为满射，$Z$ 可由拟射影 $S$ 概形表示，$R$ 可由 $S$ 概形表示，且第一投影 $R\to Z$ 可由光滑固有映射表示。则 $P$ 可由拟射影 $S$ 概形表示，且 $\alpha$ 可由光滑映射表示。
\end{origstatement}
\begin{proof}
为减轻记号，也以 $Z,R$ 表示相应的概形。由于结构映射 $Z\to S$ 拟射影，它分离；因而第一投影 $Z\times_S Z\to Z$ 也分离。但第一投影 $R\to Z$ 固有，并自然地通过映射 $h:R\to Z\times_S Z$ 分解。因此 $h$ 固有。又 $h$ 是单态射，即 $h$ 在 $T$ 点上单射。因此，由 [EGAIV3，8.11.5]，$h$ 为闭嵌入。

显然，对每个 $S$ 概形 $T$，子集 $R(T)\subset Z(T)\times_{S(T)}Z(T)$ 是某个等价关系的图像。此外，概形映射 $R\to Z$ 平坦且固有，$Z$ 是拟射影 $S$ 概形。因此存在拟射影 $S$ 概形 $Q$，以及忠实平坦射影映射 $Z\to Q$，使 $R=Z\times_Q Z$。（换言之，拟射影概形上的平坦固有等价关系是有效的。）事实上，$R$ 定义从 $Z$ 到 Hilbert 概形 $\mathbf{Hilb}_{Z/S}$ 的映射，其图像作为闭子概形位于万有概形中，并下降为 $Q\subset\mathbf{Hilb}_{Z/S}$；更多细节见 [AK80，定理 (2.9)，第 70 页]。

由于 $Z\to Q$ 平坦，若（且仅若）它的纤维光滑，则它光滑。但是，至多相差底域扩张，这些纤维与 $R\to Z$ 的纤维相同。由假设，$R\to Z$ 光滑。因此 $Z\to Q$ 光滑。

还需证明 $Q$ 表示 $P$。首先注意，$Z\to Q$ 是（表示）平展层的满射。事实上，给定 $Q(T)$ 的元素，置 $A:=Z\times_Q T$。则 $A\to T$ 光滑。所以由 [EGAIV4，17.16.3(ii)]，存在平展覆盖 $T'\to T$ 以及 $T$ 映射 $T'\to A$。于是，$T'\to A\to Z$ 为 $Z(T')$ 的元素，映到由给定 $Q(T)$ 的元素所诱导的 $Q(T')$ 的元素。
% 原书 PDF 279，印刷页 267；引理9.4.9证明续。

由于 $Z\to Q$ 为平展层的满射，根据下面的习题 9.4.10，$Q$ 在平展层范畴中是映射对 $R\rightrightarrows Z$ 的余等化子；这是一般 Grothendieck 拓扑中的一道初等习题。由同一习题，$P$ 也是这对映射的余等化子。但是，在任意范畴中，余等化子唯一到唯一同构。因此，$Q$ 表示 $P$。
\end{proof}
\begin{origstatement}{习题}{9.4.10}\label{ch09s04:9.4.10}
给定平展层的映射 $F\to G$，证明它为满射，当且仅当 $G$ 为映射对 $F\times_GF\rightrightarrows F$ 的余等化子。
\end{origstatement}
\begin{origstatement}{习题}{9.4.11}\label{ch09s04:9.4.11}
假定 $X/S$ 射影且平坦，几何纤维为整概形，$S$ 为 Noether 概形。设 $Z\subset\PS$ 为有限型子概形。证明 $Z$ 拟射影。
\end{origstatement}
\begin{origstatement}{习题}{9.4.12}\label{ch09s04:9.4.12}
假定 $f:X\to S$ 在 $S$ 上 Zariski 局部地射影，且平坦，几何纤维为整概形。首先证明，若万有层 $\PP$ 存在，则 Abel 映射 $\mathbf A_{X/S}:\mathbf{Div}_{X/S}\to\PS$ 在 $S$ 上 Zariski 局部地为射影态射。

其次证明，一般地，$\mathbf A_{X/S}:\mathbf{Div}_{X/S}\to\PS$ 为固有态射。将这一情形归结到第一种情形来证明：只需用 $f:X\to S$ 本身作基变换。
\end{origstatement}
\begin{origstatement}{习题}{9.4.13}\label{ch09s04:9.4.13}
假定 $f:X\to S$ 平坦，且在 $S$ 上 Zariski 局部地射影。假定其几何纤维 $X_k$ 为整曲线，算术亏格 $\dim H^1(\mathcal O_{X_k})$ 至少为 $1$。设 $X_0\subset X$ 为 $X/S$ 光滑的开子概形。证明存在自然闭嵌入 $A:X_0\hookrightarrow\PS$。
\end{origstatement}
\begin{origstatement}{例}{9.4.14}\label{ch09s04:9.4.14}
在定理 9.4.8 中，假定 $f$ 的几何纤维为整概形。这个条件不仅是证明成立所必需的，也是结论成立所必需的。下列归于 Mumford 的例子说明了这一点，见 [FGA，第 236-1 页] 及 [BLR90，第 210 页]。

设 $\mathbb R[[t]]$ 为形式幂级数环，$S$ 为其谱。设 $X\subset\mathbb P_S^2$ 为非齐次方程 $x^2+y^2=t$ 所定义的子概形，$f:X\to S$ 为结构映射。泛纤维 $X_\sigma$ 为非退化二次曲线。特殊纤维 $X_0$ 是一对共轭直线；$X_0$ 不可约且几何连通。因此，在引理 9.3.4 中取 $D:=X$，由蕴含关系 (iii)$\Rightarrow$(i)，$f$ 平坦。并且，由习题 9.3.11，$\mathcal O_S\xrightarrow{\sim}f_*\mathcal O_X$ 普遍成立。但是，正如将要看到的，$\PS$ 并不存在！

另一方面，置 $S':=\mathbb C[[t]]$、$X':=X\times S'$。则 $f_{S'}$ 有截面：例如，令
\[
x:=\sqrt{-1},\qquad y:=\sqrt{1+t}=1+(1/2)t-(1/8)t^2+\cdots
\]
便定义一个截面 $g'$。因此，由比较定理即定理 9.2.5，$X'/S'$ 的四个相对 Picard 函子均相等。此外，正如将要看到的，$\mathbf{Pic}_{X'/S'}$ 存在，却不分离！

事实上，$\mathbf{Pic}_{X'/S'}$ 是同构的非分离开子概形 $S'_n$（$n\in\mathbb Z$）的不交并。每个 $S'_n$ 都由 $S'$ 将原点重复无穷多次而得到；更准确地，为形成 $S'_n$，对每一对满足 $a+b=n$ 的 $a,b\in\mathbb Z$，取 $S'$ 的一个副本 $S'_{a,b}$，并将这些 $S'_{a,b}$ 在闭点之外粘合起来。每个 $S'_{a,b}$ 参数化泛纤维 $X'_\sigma$ 上次数为 $n$ 的可逆层的一种不同退化；退化后的层在第一条直线上次数为 $a$，在第二条上次数为 $b$。此外，复共轭交换这两条直线，所以也交换 $S'_{a,b}$ 与 $S'_{b,a}$。

假定 $\PS$ 存在。则由习题 9.4.4，$\PS\times S'=\mathbf{Pic}_{X'/S'}$。由于 $S'_{a,b}$ 与 $S'_{b,a}$ 的闭点互为共轭，它们映到 $\PS$ 的同一个点。此点位于某个仿射开子概形 $U$ 中。因此，这两个闭点
% 原书 PDF 280，印刷页 268；例9.4.14续。
位于 $U$ 在 $\mathbf{Pic}_{X'/S'}$ 中的逆像 $U'$ 中。但 $U'$ 仿射，因为 $U,S,S'$ 都仿射。然而，若 $a\ne b$，则这两个闭点不同，因而不在任何共同的仿射开集 $U'$ 中。矛盾。因此 $\PS$ 不可能存在。

最后证明，$\Pic_{X'/S'}$ 可由 $\coprod S'_n$ 表示。首先注意 $X'$ 正则；事实上，在 $\mathbb P_{\mathbb C}^2\times\mathbb A_{\mathbb C}^1$ 中，方程 $x^2+y^2=t$ 定义一个光滑曲面。因此，在 $X'$ 上，每条约化曲线都是有效除子。特别考虑下面三个：直线 $L:x=\sqrt{-1}y,t=0$，直线 $M:x=-\sqrt{-1}y,t=0$，以及上面定义的 $f_{S'}$ 的截面 $g$ 的像 $A$。

置 $\PP'_{a,b}:=\mathcal O_{X'}(bL+nA)$，其中 $n=a+b$。它在泛纤维上的限制 $\PP'_{a,b}|X'_\sigma$ 的次数为 $n$，因为 $L\cap X'_\sigma$ 为空，而 $A$ 为一个截面的像。又 $\PP'_{a,b}|M$ 的次数为 $b$，因为 $A\cap M$ 为空，$L\cap M$ 为一个约化 $\mathbb C$ 有理点。$\PP'_{a,b}|L$ 的次数为 $a$，因为除此之外还有 $\mathcal O_{X'}(L)\otimes\mathcal O_{X'}(M)\simeq\mathcal O_{X'}$，这是由于 $L+M$ 的理想由 $t$ 生成。最后，$S'$ 上每个可逆层均平凡；固定一个 $g$ 刚化 $\mathcal O_{S'}\xrightarrow{\sim}g^*\PP'$，并用它等同这两个层。

在 $X'\times_{S'}\coprod S'_n$ 上，通过在 $S'_{a,b}$ 上放置 $\PP'_{a,b}$，形成可逆层 $\PP'$；显然，这些 $\PP'_{a,b}$ 可拼接在一起。现在只需证明：给定任意 $S'$ 概形 $T$ 及 $X'_T$ 上任意可逆层 $\L$，存在唯一 $S'$ 映射 $q:T\to\coprod S'_n$，以及 $T$ 上某个可逆层 $\N$，使 $(1\times q)^*\PP'\simeq\L\otimes f_T^*\N$。

以 $\L\otimes f_T^*g_T^*\L^{-1}$ 替换 $\L$。则 $g_T^*\L=\mathcal O_T$，因为 $g_T^*f_T^*=1$。因此，若 $q$ 和 $\N$ 存在，则必有 $\N\simeq\mathcal O_T$，因为 $g_T^*(1\times q)^*\PP'=q^*g^*\PP'$，且 $g^*\PP'=\mathcal O_{S'}$。

显然，可以假定 $T$ 连通。则由 [EGAIII2，7.9.5]，函数 $s\mapsto\chi(X'_t,\L_t)$ 在 $T$ 上为常值。置 $n:=\chi(X'_t,\L_t)-1$。

固定满足 $a+b=n$ 的 $a,b$。置
\[
\mathcal M:=\L^{-1}\otimes(1\times\tau)^*\PP'_{a,b},\qquad
\N:=f_{T*}\mathcal M.
\]
显然，$g_T^*\mathcal M=\mathcal O_T$。形成自然映射 $u:f_T^*\N\to\mathcal M$。

设 $T_\sigma$ 为结构映射 $\tau:T\to S'$ 的泛纤维。则 $T_\sigma\subset T$ 为开集。设 $t\in T_\sigma$，则 $X'_t$ 为带有有理点 $A_t$ 的非退化平面二次曲线。所以 $X'_t\simeq\mathbb P_{k_t}^1$。因此 $\L_t\simeq\mathcal O_{X'_t}(nA_t)$，故 $\mathcal M_t\simeq\mathcal O_{X'_t}$。由 Serre 的显式计算 [EGAIII1，2.1.12]，$H^1(X'_t,\mathcal M_t)=0$，且 $H^0(X'_t,\mathcal M_t)=k_t$。

因此，由第 9.3.10 小节的理论，$\N$ 在 $t$ 处可逆，且 $\N$ 的形成与取到 $X'_t$ 上交换。所以 $u:f_T^*\N\to\mathcal M$ 的形成也与取到 $X'_t$ 上交换。而 $u$ 在 $X'_t$ 上为同构。因此由 Nakayama 引理，$u$ 沿 $X'_t$ 为满射。但 $u$ 的定义域和目标沿 $X'_t$ 均可逆。因此，$u$ 沿 $X'_t$ 为同构；由于 $t\in T_\sigma$ 任意，$u$ 在 $T_\sigma$ 上为同构。现在，$g_T^*\mathcal M=\mathcal O_T$，且 $g_T^*f_T^*=1$。所以 $\N|T_\sigma=\mathcal O_{T_\sigma}$。因此，$\L|X'_{T_\sigma}=(1\times\tau)^*\PP'_{a,b}|X'_{T_\sigma}$。

设 $T'_{a,b}$ 为满足以下条件的 $t\in T$ 的集合：$\tau(t)\in S'$ 为闭点，$\L|L_t$ 的次数为 $a$，$\L|M_t$ 的次数为 $b$。固定 $t\in T'_{a,b}$。则 $\mathcal M|L_t$ 的次数为 $0$，所以 $\mathcal M|L_t\simeq\mathcal O_{L_t}$。类似地，$\mathcal M|M_t\simeq\mathcal O_{M_t}$。考虑自然短正合列
\[
0\longrightarrow\mathcal O_{L_t}(-1)\longrightarrow\mathcal O_{X'_t}
\longrightarrow\mathcal O_{M_t}\longrightarrow0.
\]
用 $\mathcal M$ 扭转并取上同调，得 $H^1(X'_t,\mathcal M_t)=0$，且 $H^0(X'_t,\mathcal M_t)=k_t$。因此，在 $X'_t$ 上，映射 $u$ 成为映射 $\mathcal O_{X'_t}\to\mathcal M_t$。由于它限制到 $L_t$ 和 $M_t$ 后均为满射，所以这个映射为满射。因此，它为同构。

如上，可知 $u$ 在 $X'_t$ 的一个开邻域 $V'$ 上为同构。置 $W':=f_T(X'_T-V')$。由于 $f$ 固有，$W'$ 为开集。又 $f^{-1}W'\subset V'$。所以，$u$ 在 $W'$ 上为同构。因此，对所有 $t'\in W'$，有 $\mathcal O_{X'_{t'}}\xrightarrow{\sim}\mathcal M_{t'}$。所以 $W'\subset T'_{a,b}\cup T_\sigma$。
% SOURCE ISSUE：原书确实写W'=f_T(X'_T−V')，未印取补集，但随后称它开；未擅改。
% 原书 PDF 281，印刷页 269；例9.4.14续。

置 $T_{a,b}:=T'_{a,b}\cup T_\sigma$。则 $T_{a,b}\subset T$ 为开集，因为它包含其每一点的一个邻域。此外，$u$ 在 $T_{a,b}$ 上为同构。因此 $\N|T_\sigma=\mathcal O_{T_\sigma}$，仍是由于 $g_T^*\mathcal M=\mathcal O_T$ 且 $g_T^*f_T^*=1$。所以，$\L|X'_{T_{a,b}}=(1\times\tau)^*\PP'_{a,b}|X'_{T_{a,b}}$。
% SOURCE ISSUE：此处仍印T_sigma，而结论针对T_a,b；忠实保留。

令 $q_{a,b}:T_{a,b}\to\coprod S'_n$ 为结构映射 $\tau:T\to S'$ 与将 $S'$ 作为 $S'_{a,b}$ 包含到 $S'_n$ 中的包含映射的复合。显然，$(1\times q_{a,b})^*\PP'=\L|X'_{T_{a,b}}$。显然，当 $a,b,n$ 变动时，各 $q_{a,b}$ 拼接成映射 $q:T\to\coprod S'_n$，使 $(1\times q)^*\PP'=\L$。显然，这个映射 $q$ 是唯一的 $S'$ 映射，使得对 $T$ 上某个可逆层 $\N$，有 $(1\times q)^*\PP'\simeq\L\otimes f_T^*\N$。因此，$\coprod S'_n$ 表示 $\Pic_{X'/S'}$，且 $\PP'$ 是万有层。
\end{origstatement}
\begin{origstatement}{习题}{9.4.15}\label{ch09s04:9.4.15}
假定 $X=\mathbb P(\mathcal E)$，其中 $\mathcal E$ 为 $S$ 上的局部自由层，处处秩至少为 $2$。证明 $\PS$ 存在，且表示 $\Pic_{X/S}$；事实上，$\PS=\mathbb Z_S$，其中 $\mathbb Z_S$ 表示以 $\mathbb Z$ 为指标的 $S$ 的副本的不交并。
\end{origstatement}
\begin{origstatement}{习题}{9.4.16}\label{ch09s04:9.4.16}
考虑习题 9.2.4 中的曲线 $X$。证明 $\mathbf{Pic}_{X/\mathbb R}=\mathbb Z_{\mathbb R}$。
\end{origstatement}
\begin{origstatement}{命题}{9.4.17}\label{ch09s04:9.4.17}
假定 $\PS$ 存在且表示 $\PF$。则 $\PS$ 局部有限型。
\end{origstatement}
\begin{proof}
置 $P:=\PF$。由 [EGAIV3，8.14.2]，只需验证以下条件：对仿射 $S$ 概形的每个滤过逆系统 $T_i$，自然映射为双射：
\[
\varinjlim P(T_i)\xrightarrow{\sim}P(\varprojlim T_i).
\]
为验证单射性，置 $T:=\varprojlim T_i$。固定 $i$，取 $\lambda_i\in P(T_i)$。以 $X_{T'_i}$ 上的层 $\L'_i$ 表示 $\lambda_i$，其中 $T'_i\to T_i$ 为 fppf 覆盖。置 $T':=T'_i\times_{T_i}T$、$\L':=\L_i|X_{T'}$。设 $\lambda$ 为 $\lambda_i$ 在 $P(T)$ 中的像。则 $\L'$ 表示 $\lambda$。

假定 $\lambda=0\in P(T)$。则存在 fppf 覆盖 $T''\to T'$，使 $\L'|X_{T''}\simeq\mathcal O_{X_{T''}}$。由 [EGAIV3，8.8.2、8.10.5(vi)、11.2.6、8.5.2(i) 及 8.5.2.4]，存在 $j\geq i$ 及 fppf 覆盖 $T''_j\to T'_j$，其中 $T'_j:=T'_i\times_{T_i}T$，使 $\L'|X_{T''_j}\simeq\mathcal O_{X_{T''_j}}$。所以，$\lambda_i$ 映到 $0\in P(T_j)$。因此，该映射为单射。
% SOURCE ISSUE：原书此处L_i省略prime，T'_j定义末因子印T而非T_j；忠实保留。

满射性可类似证明。
\end{proof}
\begin{origstatement}{注}{9.4.18}\label{ch09s04:9.4.18}
关于 $\PS$，有三个重要的存在性定理，它们细化了定理 9.4.8。这些定理在主定理之后不久就被证明，每个都涉及新的思想。
\end{origstatement}
首先，Mumford 证明了定理 9.4.8 的下列推广，它与例 9.4.14 很好地相配合。

\begin{origstatement}{定理}{9.4.18.1}\label{ch09s04:9.4.18.1}
假定 $X/S$ 射影且平坦，几何纤维约化且连通；假定其普通纤维的不可约分支几何不可约。则 $\PS$ 存在。
\end{origstatement}
Mumford 在 [Mum66] 第 viii 页下部陈述了这个定理。他说，证明是 [Mum66，第 19--21 讲] 中证明的推广；那里考虑 $S$ 为代数闭域的谱、$X$ 为光滑曲面的情形。那个证明基于他的独立零维圈理论。[AK79，第 23--28 页] 进一步发展了这一理论，并用它证明了 [AK79，定理 (3.1)]；该定理断言，当 $X/S$ 平坦且局部射影，几何纤维为整概形时，$\PS$ 有自然紧化。

另一方面，Grothendieck [FGA，第 236-1 页] 将一个略有不同的定理归于 Mumford；在那个版本中，既不假定几何纤维连通，也不假定普通
% 原书 PDF 282，印刷页 270；上段续。
纤维连通（另见 [BLR90，第 210 页]）。Grothendieck 说，证明基于 Mumford 商构造的一个细化，并引述了即将出版的、Mumford 和 Tate 于 1962 年春在 Harvard 举行的讨论班的讲义。

2003 年 11 月，Mumford 好意将他个人保存的该讨论班手写笔记夹寄给本章作者；文件夹标有“Groth-Mumford-Tate”，包含三人各自报告的笔记，以及各自撰写的笔记。几乎全部内容都已在别处发表；Mumford 的部分见于他的著作 [Mum65]、[Mum66] 及 [Mum70]。

笔记中有 Mumford 亲笔写下的定理的精确陈述及证明粗略梗概。这一陈述也与定理 9.4.18.1 略有不同：他划去了几何纤维连通的假设，改用较弱的普通纤维连通的假设。证明总体上类似于他在 [Mum66，第 19--21 讲] 中的证明。

其次，Grothendieck 证明了下列“泛可表性”定理。

\begin{origstatement}{定理}{9.4.18.2}\label{ch09s04:9.4.18.2}
假定 $X/S$ 固有，$S$ 为整概形。则存在非空开子集 $V\subset S$，使得 $\mathbf{Pic}_{X_V/V}$ 存在，表示 $\Pic_{(X_V/V)(\mathrm{fppf})}$，并为拟射影开子概形的不交并。
\end{origstatement}
一个特别重要的特殊情形包含在下面的推论中。

\begin{origstatement}{推论}{9.4.18.3}\label{ch09s04:9.4.18.3}
假定 $S$ 为域 $k$ 的谱，$X$ 完备。则 $\mathbf{Pic}_{X/k}$ 存在，表示 $\Pic_{(X/k)(\mathrm{fppf})}$，并为拟射影开子概形的不交并。
\end{origstatement}
在 Grothendieck 发现定理 9.4.18.2 之前，他在假定 $X/k$ 射影的情形下证明了推论 9.4.18.3 [FGA，推论 6.6，第 232-17 页]。为此，他发展了“相对可表性”的方法，由定理 9.4.8 推出其他情形下 Picard 概形的存在性。这个方法结合了 Oort 的一种“逐层分解”（dévissage）[Oor62，§6]；后者将 Picard 概形给出为 $\mathbf{Pic}_{X_{\mathrm{red}}/k}$ 的一个群子概形被仿射群概形的扩张。

在 [FGA，注 6.6，第 232-17 页] 中，Grothendieck 说，推论 9.4.18.3 在一般情形成立是“极其可信的”，并且可通过推广相对可表性的方法来证明，使它涵盖 $X'$ 射影的满射 $X'\to X$ 的情形，例如 Chow 引理 [EGAII，定理 5.6.1] 所给出的映射。此外，他在 [FGA，注 6.6，第 232-17 页] 中猜想，对域上固有概形之间的任意满射 $X'\to X$，所诱导的 Picard 概形之间的映射是仿射的。他说，这个猜想源自对“非平坦下降”的考虑；这是一种不要求映射平坦、而要求对象平坦的下降理论。

Murre [Mur64，(II.15)] 在感谢 Grothendieck 的帮助的同时，给出了推论 9.4.18.3 核心部分的第一个证明：若 $X/k$ 固有，则 $\mathbf{Pic}_{X/k}$ 存在且局部有限型。Murre 没有使用相对可表性的方法，而是找出了七个条件 [Mur64，(1.2.1)]，它们是从域上概形到 Abel 群的函子可表的充要条件。随后，他证明，在 fpqc 拓扑中局部化的 Picard 函子满足这七个条件。为处理最后两个条件，他使用了 Chow 引理和定理 9.4.8。

与此同时，据 Murre [Mur64，第 5 页]，Grothendieck 已经证明了定理 9.4.18.2。后来，证明分两部分发表。第一部分
\endgroup
